% Mapa navegable de problemas por categoría.
% Este archivo se inserta antes del índice desde Problemset1.tex.

\definecolor{appheader}{HTML}{0B6E99}
\definecolor{approw}{HTML}{E8F7FC}
\definecolor{theoryheader}{HTML}{6D44A5}
\definecolor{theoryrow}{HTML}{F2ECFA}
\definecolor{compheader}{HTML}{087F5B}
\definecolor{comprow}{HTML}{E7F8F1}

\definecolor{diffone}{HTML}{D9F7BE}
\definecolor{difftwo}{HTML}{B7EB8F}
\definecolor{diffthree}{HTML}{95DE64}
\definecolor{difffour}{HTML}{D3F261}
\definecolor{difffive}{HTML}{FFEC3D}
\definecolor{diffsix}{HTML}{FFC53D}
\definecolor{diffseven}{HTML}{FA8C16}
\definecolor{diffeight}{HTML}{FF7A45}
\definecolor{diffnine}{HTML}{FF4D4F}
\definecolor{difficult}{HTML}{CF1322}

\newcommand{\dificultad}[1]{%
  \begingroup
  \ifcase#1\or
    \cellcolor{diffone}\or
    \cellcolor{difftwo}\or
    \cellcolor{diffthree}\or
    \cellcolor{difffour}\or
    \cellcolor{difffive}\or
    \cellcolor{diffsix}\or
    \cellcolor{diffseven}\or
    \cellcolor{diffeight}\or
    \cellcolor{diffnine}\color{white}\or
    \cellcolor{difficult}\color{white}%
  \fi
  \centering\bfseries #1/10
  \endgroup
}

\newcommand{\enlaceproblema}[2]{%
  \hyperref[#1]{\textbf{\ref*{#1}}} &
  \hyperref[#1]{#2}%
}
\newcommand{\filaaplicacion}[4]{%
  \rowcolor{approw}\enlaceproblema{#1}{#2} & #3 & \dificultad{#4}\\\hline
}
\newcommand{\filateorica}[4]{%
  \rowcolor{theoryrow}\enlaceproblema{#1}{#2} & #3 & \dificultad{#4}\\\hline
}
\newcommand{\filacomputacional}[4]{%
  \rowcolor{comprow}\enlaceproblema{#1}{#2} & #3 & \dificultad{#4}\\\hline
}

\newcolumntype{N}{>{\centering\arraybackslash}p{1.0cm}}
\newcolumntype{T}{>{\raggedright\arraybackslash}p{4.2cm}}
\newcolumntype{P}{>{\raggedright\arraybackslash}p{7.6cm}}
\newcolumntype{D}{>{\centering\arraybackslash}p{1.5cm}}

\section*{Guía de problemas por categoría}
\phantomsection
\addcontentsline{toc}{section}{Guía de problemas por categoría}

Las tablas siguientes permiten navegar directamente a cada problema. Un mismo
problema puede aparecer en más de una categoría cuando combina análisis,
aplicación científica y trabajo computacional. La dificultad considera la carga
conceptual, la extensión del desarrollo y la complejidad de implementación.

\begin{center}
\small
\begin{tabular}{c@{\hspace{0.25cm}}l@{\hspace{0.7cm}}c@{\hspace{0.25cm}}l@{\hspace{0.7cm}}c@{\hspace{0.25cm}}l}
\cellcolor{diffone}\rule{0pt}{2.2ex}\hspace{0.7cm} & 1--2: inicial &
\cellcolor{difffour}\hspace{0.7cm} & 3--4: básica &
\cellcolor{diffsix}\hspace{0.7cm} & 5--6: intermedia \\
\cellcolor{diffeight}\rule{0pt}{2.2ex}\hspace{0.7cm} & 7--8: avanzada &
\cellcolor{difficult}\hspace{0.7cm} & 9--10: muy avanzada &
& \\
\end{tabular}
\end{center}

\begin{tcolorbox}[
  colback=appheader!8,
  colframe=appheader,
  boxrule=0.8pt,
  arc=2mm,
  left=2mm,right=2mm,top=1mm,bottom=1mm]
  \centering\large\bfseries\color{appheader} Aplicación
\end{tcolorbox}

\arrayrulecolor{appheader!45}
\renewcommand{\arraystretch}{1.25}
\begin{longtable}{NTPD}
\rowcolor{appheader}
\color{white}\textbf{Sec.} &
\color{white}\textbf{Problema} &
\color{white}\textbf{Tópicos principales} &
\color{white}\textbf{Dif.}\\\hline
\endfirsthead
\multicolumn{4}{l}{\small\itshape Aplicación (continuación)}\\
\rowcolor{appheader}
\color{white}\textbf{Sec.} &
\color{white}\textbf{Problema} &
\color{white}\textbf{Tópicos principales} &
\color{white}\textbf{Dif.}\\\hline
\endhead
\filaaplicacion{sec:problema-02}{Clasificación morfológica de galaxias}{Espacios muestrales finitos; eventos; unión e intersección; secuencias.}{3}
\filaaplicacion{sec:problema-03}{Seguimiento de cometas}{Espacio infinito discreto; tiempo de espera; distribución geométrica.}{4}
\filaaplicacion{sec:problema-04}{Coordenadas ecuatoriales}{Espacios continuos; densidad uniforme; normalización; probabilidad por áreas.}{5}
\filaaplicacion{sec:problema-05}{Taxonomía de asteroides}{Espacio muestral; eventos compuestos; unión, intersección y complementos.}{3}
\filaaplicacion{sec:problema-06}{Modelado gravitacional}{Secuencias binarias; cardinalidad; eventos y estrategias de navegación.}{5}
\filaaplicacion{sec:problema-11}{Señales de ondas gravitacionales}{Clasificación; probabilidades a priori; eventos exhaustivos y falsos positivos.}{4}
\filaaplicacion{sec:problema-12}{Modos de falla de un rover}{Unión de fallas; complemento; confiabilidad y ciclos redundantes.}{4}
\filaaplicacion{sec:problema-13}{Albedo radar de asteroides}{Probabilidad total; Bayes; sensibilidad de clasificación.}{5}
\filaaplicacion{sec:problema-14}{Pipeline de transitorios}{Regla de multiplicación; eventos dependientes; etapas secuenciales.}{4}
\filaaplicacion{sec:problema-15}{Confiabilidad de una sonda}{Independencia; fallas exactas; confiabilidad de sistemas.}{6}
\filaaplicacion{sec:problema-16}{Interferometría gravitacional}{Independencia; coincidencias; ruido instrumental; detección.}{6}
\filaaplicacion{sec:problema-21}{Clasificación mineralógica}{Inferencia bayesiana; verosimilitud; espectroscopía; decisión científica.}{5}
\filaaplicacion{sec:problema-22}{Aislamiento de fallas ADCS}{Diagnóstico bayesiano; causas raíz; telemetría; probabilidad posterior.}{5}
\filaaplicacion{sec:problema-23}{Clasificación orbital de cometas}{Bayes; errores de medición; hipótesis parabólica e hiperbólica.}{6}
\filaaplicacion{sec:problema-24}{Carga útil de una sonda}{Principio multiplicativo; producto cartesiano; configuraciones restringidas.}{4}
\filaaplicacion{sec:problema-26}{Observación con telescopio robótico}{Permutaciones; posiciones prioritarias; restricciones de adyacencia.}{6}
\filaaplicacion{sec:problema-27}{Apuntamiento de telescopio espacial}{Permutaciones sin repetición; restricciones térmicas; conteo complementario.}{5}
\filaaplicacion{sec:problema-28}{Espectrómetro de masas}{Secuencias restringidas; permutaciones; backtracking; redes.}{7}
\filaaplicacion{sec:problema-29}{Sitios de aterrizaje}{Combinaciones; selección estratificada; probabilidad; Monte Carlo.}{6}
\filaaplicacion{sec:problema-30}{Sub-arrays interferométricos}{Combinaciones; líneas de base; restricciones; optimización y geometría.}{8}
\end{longtable}

\clearpage
\begin{tcolorbox}[
  colback=theoryheader!8,
  colframe=theoryheader,
  boxrule=0.8pt,
  arc=2mm,
  left=2mm,right=2mm,top=1mm,bottom=1mm]
  \centering\large\bfseries\color{theoryheader} Teórico
\end{tcolorbox}

\arrayrulecolor{theoryheader!45}
\begin{longtable}{NTPD}
\rowcolor{theoryheader}
\color{white}\textbf{Sec.} &
\color{white}\textbf{Problema} &
\color{white}\textbf{Tópicos principales} &
\color{white}\textbf{Dif.}\\\hline
\endfirsthead
\multicolumn{4}{l}{\small\itshape Teórico (continuación)}\\
\rowcolor{theoryheader}
\color{white}\textbf{Sec.} &
\color{white}\textbf{Problema} &
\color{white}\textbf{Tópicos principales} &
\color{white}\textbf{Dif.}\\\hline
\endhead
\filateorica{sec:problema-01}{Identificación de espacios muestrales}{Resultados ordenados; defectos; conteos; extracción con y sin reemplazo.}{2}
\filateorica{sec:problema-07}{Probabilidad del evento nulo}{Axiomas de probabilidad; aditividad numerable; demostración de $P(\emptyset)=0$.}{3}
\filateorica{sec:problema-08}{Inclusión-exclusión}{Uniones e intersecciones; sumas alternantes; inducción matemática.}{7}
\filateorica{sec:problema-09}{Probabilidad I}{Axiomas; distribuciones de Poisson y geométrica; contraejemplos.}{6}
\filateorica{sec:problema-10}{Probabilidad II}{Medidas sobre Borel; densidades; normalización; contraejemplos.}{7}
\filateorica{sec:problema-17}{Probabilidad condicional I}{Cotas probabilísticas; inclusión-exclusión; demostración algebraica.}{6}
\filateorica{sec:problema-18}{Probabilidad condicional II.1}{Distribución del tamaño familiar; mezcla binomial; normalización.}{7}
\filateorica{sec:problema-19}{Probabilidad condicional II.2}{Distribuciones familiares; conteo de varones; condicionamiento.}{8}
\filateorica{sec:problema-20}{Actualización secuencial}{Teorema de Bayes; posterior secuencial; independencia condicional.}{7}
\filateorica{sec:problema-25}{Urna y canicas}{Extracción sin reemplazo; primera ocurrencia; productos; límite geométrico.}{8}
\filateorica{sec:problema-31}{Independencia de eventos}{Distribución exponencial; supervivencia; complemento; mínimo de vidas.}{5}
\end{longtable}

\clearpage
\begin{tcolorbox}[
  colback=compheader!8,
  colframe=compheader,
  boxrule=0.8pt,
  arc=2mm,
  left=2mm,right=2mm,top=1mm,bottom=1mm]
  \centering\large\bfseries\color{compheader} Computacional
\end{tcolorbox}

\arrayrulecolor{compheader!45}
\begin{longtable}{NTPD}
\rowcolor{compheader}
\color{white}\textbf{Sec.} &
\color{white}\textbf{Problema} &
\color{white}\textbf{Tópicos y herramientas} &
\color{white}\textbf{Dif.}\\\hline
\endfirsthead
\multicolumn{4}{l}{\small\itshape Computacional (continuación)}\\
\rowcolor{compheader}
\color{white}\textbf{Sec.} &
\color{white}\textbf{Problema} &
\color{white}\textbf{Tópicos y herramientas} &
\color{white}\textbf{Dif.}\\\hline
\endhead
\filacomputacional{sec:problema-02}{Clasificación morfológica}{Producto cartesiano; \texttt{itertools}; enumeración exhaustiva.}{3}
\filacomputacional{sec:problema-03}{Seguimiento de cometas}{Simulación geométrica; tiempos de espera; histograma de frecuencias.}{4}
\filacomputacional{sec:problema-04}{Coordenadas ecuatoriales}{Muestreo uniforme; dispersión 2D; estimación Monte Carlo.}{5}
\filacomputacional{sec:problema-07}{Evento nulo}{Sumas parciales; divergencia; gráfica y tabla comparativa.}{3}
\filacomputacional{sec:problema-11}{Señales gravitacionales}{\texttt{np.random.choice}; frecuencias relativas; validación empírica.}{4}
\filacomputacional{sec:problema-14}{Pipeline de transitorios}{Simulación secuencial; gráfico de embudo; análisis de sensibilidad.}{4}
\filacomputacional{sec:problema-15}{Confiabilidad de sonda}{Simulación matricial; conteo de fallas; análisis de sensibilidad.}{6}
\filacomputacional{sec:problema-16}{Interferometría gravitacional}{Eventos raros; tasa de falsas alarmas; gráfica de convergencia.}{6}
\filacomputacional{sec:problema-21}{Clasificación mineralógica}{Simulación bayesiana; diagrama de Sankey con \texttt{plotly}; sensibilidad.}{6}
\filacomputacional{sec:problema-23}{Órbitas de cometas}{Simulación de clasificación; matriz de confusión; sensibilidad instrumental.}{6}
\filacomputacional{sec:problema-24}{Configuración de carga útil}{\texttt{itertools.product}; filtrado; espacio de configuraciones.}{4}
\filacomputacional{sec:problema-26}{Telescopio robótico}{Factoriales; \texttt{itertools}; filtrado de adyacencias; Monte Carlo.}{6}
\filacomputacional{sec:problema-28}{Espectrómetro de masas}{Permutaciones; Monte Carlo; backtracking; \texttt{networkx}.}{7}
\filacomputacional{sec:problema-29}{Sitios de aterrizaje}{\texttt{combinations}; Monte Carlo; gráfica de convergencia.}{6}
\filacomputacional{sec:problema-30}{Sub-arrays interferométricos}{Combinaciones; geometría aleatoria; \texttt{seaborn}; fuerza bruta.}{8}
\end{longtable}

\arrayrulecolor{black}
\renewcommand{\arraystretch}{1}
